\documentclass[a4paper,12pt]{article}

\usepackage[ngerman]{babel}
\usepackage[utf8]{inputenc}
\usepackage[T1]{fontenc}
\usepackage{lmodern}
\usepackage{geometry}
\usepackage{listings}
\usepackage{xcolor}

\geometry{margin=2.5cm}

\title{Praktische Übung: Schnittstellen für eine Webanwendung entwerfen}
\author{Modul 321}
\date{\today}

\lstset{
    basicstyle=\ttfamily\small,
    breaklines=true,
    frame=single,
    backgroundcolor=\color{gray!10}
}

\begin{document}

\maketitle

\section*{Ziel der Übung}

In dieser Übung entwerft und implementiert ihr Schnittstellen für eine einfache Anwendung.
Ziel ist nicht eine vollständige Produktentwicklung, sondern das saubere Einordnen
und Anwenden von Protokollen und Schnittstellenstilen.

\section*{Ausgangslage}

Ihr entwickelt eine kleine Bestell-Anwendung mit diesen Anforderungen:

\begin{itemize}
\item Benutzer melden sich an und verwalten ihre Bestellungen
\item Eine mobile App soll nur die Daten laden, die sie aktuell braucht
\item Interne Dienste (Lager und Bezahlung) müssen miteinander kommunizieren
\item Benutzer sollen sofort benachrichtigt werden, wenn ihre Bestellung versendet wurde
\item Alle externen Anfragen laufen über einen einzigen Zugangspunkt
\end{itemize}

\section*{Variante A: Konzeptioneller Entwurf}

Erstellt ein Dokument oder ein Diagramm, das zeigt:

\begin{itemize}
\item Welche Schnittstellenstile für welchen Bereich eingesetzt werden
\item Wie Client, Gateway und interne Dienste zusammenarbeiten
\item Welche Protokolle verwendet werden
\item Kurze Begründung für jede Wahl
\end{itemize}

\section*{Variante B: Technische Demo}

Implementiert eine kleine Demo mit mindestens \textbf{zwei} der folgenden Bausteine:

\begin{lstlisting}
bestell-app/
|
|-- gateway/
|   `-- index.js        # Einfaches Routing zu den Diensten
|
|-- api-rest/
|   `-- bestellungen.js # REST-Endpunkte fuer Bestellungen
|
|-- api-graphql/
|   `-- schema.js       # GraphQL-Schema fuer Benutzerdaten
|
|-- intern-grpc/
|   `-- lager.proto     # Protobuf-Definition fuer Lagerdienst
|
|-- ws-live/
|   `-- versand.js      # WebSocket fuer Versandbenachrichtigung
|
`-- app.js
\end{lstlisting}

Die Implementierung darf sehr einfach bleiben. Es reicht, wenn klar wird, welche
Verantwortung jede Schnittstelle übernimmt.

\section*{Arbeitsaufträge}

\subsection*{Aufgabe 1: Schnittstellen benennen und begründen}
Nennen Sie für jeden der folgenden Bereiche den passenden Schnittstellenstil:
\begin{enumerate}
\item Benutzer-Login und Bestellverwaltung (extern, öffentlich)
\item Mobile App, die je nach Ansicht unterschiedliche Daten braucht
\item Kommunikation zwischen Lager- und Bezahldienst (intern, performant)
\item Versandbenachrichtigung in Echtzeit
\end{enumerate}
Begründen Sie jede Wahl in einem Satz.

\subsection*{Aufgabe 2: Gateway beschreiben}
Beschreiben Sie, was das Gateway in dieser Anwendung übernimmt:
\begin{enumerate}
\item Welche Anfragen leitet es wohin weiter?
\item Welche Querschnittsaufgaben übernimmt es zentral?
\item Was würde passieren, wenn das Gateway fehlt?
\end{enumerate}

\subsection*{Aufgabe 3: REST-Endpunkte entwerfen}
Entwerfen Sie die wichtigsten REST-Endpunkte für die Bestellverwaltung:
\begin{itemize}
\item Alle Bestellungen eines Benutzers abrufen
\item Eine neue Bestellung erstellen
\item Eine Bestellung löschen
\end{itemize}
Geben Sie jeweils Methode, URL und den erwarteten Status-Code bei Erfolg an.

\subsection*{Aufgabe 4: Systemfluss beschreiben}
Beschreiben Sie kurz den Ablauf, wenn ein Benutzer eine Bestellung aufgibt:
\begin{enumerate}
\item Anfrage gelangt an das Gateway
\item Authentifizierung wird geprüft
\item Bestellung wird angelegt
\item Lagerdienst und Bezahldienst werden intern informiert
\item Benutzer erhält eine Benachrichtigung
\end{enumerate}
Welcher Schnittstellenstil wird in jedem Schritt verwendet?

\subsection*{Aufgabe 5: Minimalarchitektur}
Zeichnen oder beschreiben Sie eine einfache Zielarchitektur.
Zeigen Sie mindestens:
\begin{itemize}
\item Wo das Gateway steht
\item Welche externen Schnittstellen Clients nutzen
\item Wie interne Dienste kommunizieren
\item Wo WebSockets eingesetzt werden
\end{itemize}

\section*{Abgabe}

Die Abgabe soll enthalten:

\begin{itemize}
\item Kurze Beschreibung der Anwendung
\item Übersicht über die gewählten Schnittstellen und deren Begründung
\item Skizze, Diagramm oder Code-Demo
\item REST-Endpunkt-Entwurf aus Aufgabe 3
\item Kurze Reflexion über Vor- und Nachteile der gewählten Stile
\end{itemize}

\section*{Reflexionsfragen}

\begin{enumerate}
\item Warum ist es sinnvoll, für interne und externe Kommunikation unterschiedliche Stile zu wählen?
\item Was wäre problematisch, wenn die mobile App denselben Stil wie die internen Dienste nutzen würde?
\item Wo könnte zu viel Komplexität für ein kleines Projekt eher unnötig sein?
\item Welcher Schnittstellenstil wäre am einfachsten zu ersetzen, ohne andere Teile des Systems zu ändern?
\end{enumerate}

\section*{Bewertung}

\begin{itemize}
\item Schnittstellenstile sind korrekt erkannt und dem Anwendungsfall zugeordnet
\item Die Begründungen sind nachvollziehbar und fachlich sauber
\item Der Gateway-Einsatz ist verstanden und beschrieben
\item Der REST-Entwurf ist konsistent (Methode, URL, Status-Code)
\item Die Lösung bleibt angemessen einfach für den Umfang der Aufgabe
\end{itemize}

\end{document}
